\documentclass[10pt]{article}

\usepackage[utf8]{inputenc}

\usepackage[T1]{fontenc}

\usepackage{amsmath}

\usepackage{amsfonts}

\usepackage{amssymb}

\usepackage[version=4]{mhchem}

\usepackage{stmaryrd}

\usepackage{bbold}



\begin{document}

БлОХА УгОЛ - угол отклонения $\Theta(t)$ классического Блоха вектора $\boldsymbol{R}(t)$ от оси $(-z)$ в сферической системе координат. Б. у. определяется равенствами



\[

\begin{gathered}

R^{x}(t)=\rho(t) \sin \Theta(t) \cos \varphi(t), \\

R^{y}(t)=\rho(t) \sin \Theta(t) \sin \varphi(t), \\

R^{z}(t)=-\rho(t) \cos \Theta(t) .

\end{gathered}

\]



Функции $\rho(t), \Theta(t)$ и $\varphi(t)$ находятся из уравнений, получающихся при подстановке данного представления вектора Блоха в уравнения Блоха. БЛОХА УРАВНЕНИЕ - см. Блоха вектор, Плотности матрица.



БЛОХА УСЛОВИЕ - см. Блоха теорема.



БлОХА ФУНКЦИЯ - см. Блоха теорема.



БЛОХА - МАКСВЕЛЛА УРАВНЕНИЯ - основная система уравнений полуклассической теории взаимодействия электромагнитного излучения с резонансной двухуровневой системой атомов (молекул, спинов и т. п.).



Как известно, алгебра двухуровневых излучателей является алгеброй матриц Паули $\sigma_{1}, \sigma_{2}, \sigma_{3}$ :



\[

\sigma_{1}=\left\|\begin{array}{ll}

0 & 1 \\

1 & 0

\end{array}\right\|, \quad \sigma_{2}=\left\|\begin{array}{cc}

0 & -i \\

i & 0

\end{array}\right\|, \quad \sigma_{3}=\left\|\begin{array}{cc}

1 & 0 \\

0 & -1

\end{array}\right\| .

\]



Квантовомеханич. среднее $\left\langle\sigma_{3}\right\rangle$ имеет смысл разности населенностей уровней, а средние $\left\langle\sigma_{1}\right\rangle$, $\left\langle\sigma_{2}\right\rangle$, равные нулю в чистом состоянии, в суперпозиционном состоянии имеют смысл компонент тока перехода. Удобнее вместо $\sigma_{1}$ и $\sigma_{2}$ использовать $\sigma_{ \pm}=1 / 2\left(\sigma_{1} \pm i \sigma_{2}\right)$. Они связаны следующим образом с операторами плотности тока атомного перехода



\[

\hat{j}(r, t)=\sum_{i=1}^{N}\left\{\left\langle+\left|\hat{j}_{0}^{+}\left(r-r_{i}\right)\right|-\right\rangle \sigma_{i^{+}}(t)+\left\langle-\left|\hat{j}_{0}^{-}\left(r-r_{i}\right)\right|+\right\rangle \sigma_{i^{-}}(t)\right\},

\]



где $N$ - число атомов в системе, $|+\rangle,|-\rangle$ - волновые функции возбужденного и основного состояний атома, $\hat{\boldsymbol{j}}_{0}^{+}(\boldsymbol{r})-$ положительно-частотная часть оператора плотности тока перехода отдельного атома в модели точечных атомов



\[

\left\langle+\left|\hat{j}_{0}^{+}\left(r-\boldsymbol{r}_{i}\right)\right|-\right\rangle=\boldsymbol{m}_{i} \delta\left(\boldsymbol{r}-\boldsymbol{r}_{i}\right),

\]



где $\boldsymbol{m}_{i}-$ матричный элемент тока перехода.



Система Б.- М.у. для векторного потенциала $A(r, t)$ и средних $\left\langle\sigma_{ \pm}\right\rangle,\left\langle\sigma_{3}\right\rangle$ имеет вид



\[

\begin{aligned}

& \operatorname{rot} \operatorname{rot} \boldsymbol{A}(\boldsymbol{r}, t)+\frac{1}{c^{2}} \frac{\partial^{2}}{\partial t^{2}} \boldsymbol{A}(\boldsymbol{r}, t)=\frac{4 \pi}{c} \boldsymbol{j}(\boldsymbol{r}, t), \\

& \frac{\partial\left\langle\sigma_{i^{+}}\right\rangle}{\partial t}=i \omega_{0 i}\left\langle\sigma_{i^{+}}\right\rangle+\frac{i}{\hbar c} \boldsymbol{m}_{i}^{*} A\left(\boldsymbol{r}_{i}, t\right)\left\langle\sigma_{i 3}\right\rangle, \\

& \frac{\partial\left\langle\sigma_{i 3}\right\rangle}{\partial t}=\frac{2 i}{\hbar c}\left(\boldsymbol{m}_{i}\left\langle\sigma_{i^{+}}\right\rangle-\boldsymbol{m}_{i}^{*}\left\langle\sigma_{i^{-}}\right\rangle\right) \boldsymbol{A}\left(\boldsymbol{r}_{i}, t\right),

\end{aligned}

\]



где $\hbar \omega_{0 i}$ - энергия возбуждения $i$-го атома.



Часто используется приближение медленно меняющихся амплитуд, состоящее в пренебрежении вторыми производными от медленно меняющихся амплитуд $\boldsymbol{A}_{\boldsymbol{k}}^{ \pm}(\boldsymbol{r}, t)$, определяющихся выражением



\[

A(\boldsymbol{r}, t)=\sum_{k} A_{k}^{+}(\boldsymbol{r}, t) e^{i\left(\omega_{k} t-k r\right)}+A_{k}^{-}(\boldsymbol{r}, t) e^{-i\left(\omega_{k} t-k r\right)} .

\]



Тогда первое уравнение системы принимает вид



где



\[

\left(\frac{\partial}{\partial t}+c \frac{\partial}{\partial x}\right) A_{k}^{+}(\boldsymbol{r}, t)=-i \frac{2 \pi}{k} \boldsymbol{m} \boldsymbol{R}_{\boldsymbol{k}}^{+},

\]



\[

R_{k}^{+}=\frac{1}{V_{1}} \sum_{i \in V_{1}(\boldsymbol{r})}\left\langle\sigma_{i^{+}}\right\rangle e^{i k r_{i}} e^{-i \omega_{k} t}

\]



и положено, что в атомной системе снято вырождение, так что все $\boldsymbol{m}_{i}=\boldsymbol{m}$.



А. В. Андреев. БЛУХДАЮЩАЯ ТОЧКА - см. Неблуждающая точка. БЛЯШКЕ МНОХИТЕЛЬ - см. Бляшке произведение. БЛЯШКЕ ПРОИЗВЕДЕНИЕ, Бл яшКе Функция,регулярная аналитическая функция комплексного переменного $z$, определенная в единичном круге $K=\{z ;|z|<1\}$ в виде конечного или бесконечного произведения



\[

B(z)=z^{n} \prod_{k} \frac{\left|a_{k}\right|}{a_{k}} \frac{a_{k}-z}{1-\bar{a}_{k} z},

\]



где $n$ - целое неотрицательное число, $\left\{a_{k}\right\}, k=1,2, \ldots,-$ последовательность точек $a_{k} \in K \backslash\{0\}$ такая, что произведение в правой части $\left({ }^{*}\right)$ сходится (условие сходимости необходимо лишь в случае бесконечного произведения). Б. п. было введено В. Бляшке (W. Blaschke, 1915), установившим следующее утверждение: последовательность $\left\{a_{k}\right\}$ точек $a_{k} \in K \backslash\{0\}$ определяет функцию вида $\left(^{*}\right)$ тогда и только тогда, когда сходится ряд



\[

\sum_{k}\left(1-\left|a_{k}\right|\right)

\]



Каждый множитель вида



\[

b_{k}(z)=\frac{\left|a_{k}\right|}{a_{k}} \frac{a_{k}-z}{1-\bar{a}_{k} z},

\]



называемый множи телем Бляшке для $a_{k}$, осуществляет однолистное конформное отображение круга $K$ на себя, переводящее точку $z=a_{k}$ в нуль, с нормировкой $b_{k}\left(-a_{k} /\left|a_{k}\right|\right)=1$. Множители вида $b_{0}(z)=z$ можно интерпретировать как множители Бляшке, соответствующие нулю $z=0$ и нормировке $b_{0}(1)=1$. Определение множителей Бляшке и Б. п. легко переносится на круг произвольного радиуса, а также на любую односвязную область, конформно эквивалентную кругу.



П. М. Тамразов.



БлЯШКЕ ФУНКЦИЯ - см. Бляшке произведение.



БОБЫЛЕВА-СТЕКЛОВА СЛУЧАЙ - см. Твердое тело; движение около неподвижной точки.



БОГОЛЮБОВА АКСИОМАТИЧЕСКИЙ ПОДХОД в квантовой теори и поля- см. Аксиоматическая квантовая теория поля.



БОГОЛЮБОВА ДУАЛЬНЫЕ УРАВНЕНИЯ - ЦеПОчКа эволюционных уравнений вида



\[

\begin{gathered}

\frac{\partial}{\partial t} f_{s}\left(t, x_{1}, \ldots, x_{s}\right)=\left[f_{s}\left(t, x_{1}, \ldots, x_{s}\right), H_{s}\right]- \\

\dot{v} \\

-\sum_{i \neq j=1}^{s} \frac{\partial}{\partial q_{i}} \Phi\left(q_{i}-q_{j}\right) \frac{\partial}{\partial p_{i}} f_{s-1}\left(t, x_{1}, \ldots, x_{s}\right), \quad s=1,2, \ldots,

\end{gathered}

\]



где $x_{i} \equiv\left(q_{i}, p_{i}\right) \in \mathbb{R}^{v} \times \mathbb{R}^{v}, H_{s}$ - гамильтониан системы частиц, взаимодействующих посредством потенциала $\Phi(q),[\cdot, \cdot]$ скобка Пуассона. Б. Д. У. являются сопряженными к Боголюбова уравнениям. Решение задачи Коши для Б. д.у. представляется формулой (см. [1]):



\[

\begin{aligned}

& f_{s}\left(t, x_{1}, \ldots, x_{s}\right)=\sum_{n=0}^{s} \sum_{k=0}^{n} \frac{(-1)^{k}}{k !(n-k) !} \times \\

& i_{1} i_{k} \quad i_{1} \quad i_{n} \\

& \times \sum^{*} S^{s-k}\left(t, x_{1}, \ldots, x_{s}\right) f_{s-n}^{0}\left(x_{1}, \ldots \vee, x_{s}\right) \text {, }

\end{aligned}

\]



где $S^{s-k}(t)$ - эволюционный оператор уравнения Лиувилля, а суммирование $\sum^{*}$ проводится от 1 до $s$ по $i_{1}, \ldots, i_{n}=1, i_{1} \neq i_{n}$.



Лит.: [1] ПетринаД.Я., Герасименко В. И., Малыше в П. В., Математические основы классической статистической механики, Киев, 1985.



В. И. Герасименко.



БОГОЛЮБОВА КАНОНИЧЕСКИЕ ПРЕОБРАЗОВАния - линейные преобразования операторов уничтожения и рождения частиц к операторам уничтожения и рождения квазичастиц для неидеальных ферми- и бозе-газов. Предложены Н. Н. Боголюбовым в 1947 для бозе-газа и в 1958 для ферми-газа.



Для неидеального бозе-газа Б. к. п. таковы:



\[

\begin{aligned}

& b_{k}=u_{k} \xi_{k}+v_{k} \xi_{-k}^{+}, \\

& b_{k}^{+}=u_{k} \xi_{k}^{+}+v_{k} \xi_{-k},

\end{aligned}

\]



где $b_{\boldsymbol{k}}, b_{\boldsymbol{k}}^{+}$- операторы уничтожения и рождения частиц в состоянии с импульсом $\boldsymbol{k}$, удовлетворяющие перестановоч-





\end{document}